% =====================================================================
%  1.10 Gestão ambiental e social
% =====================================================================
\section{Gestão ambiental e social}
\label{sec:salvaguardas}

As medidas ambientais e sociais acompanham todas as fases da ação
(\autoref{fig:socioambiental}). O \autoref{qua:socioambiental} mostra, por
tema, como a ação trata cada risco e como isso será comprovado
\cite[Quadros 25 e 26]{mop2026}.

\begin{figure}[htbp]
  \centering
  \caption{Medidas ambientais e sociais ao longo da ação}
  \label{fig:socioambiental}
  \fluxograma{fluxograma-socioambiental}
  \fonte{Elaborado pelo IAT (2026) com base em \citeonline[Quadros 25 e 26]{mop2026}.}
\end{figure}

\begin{quadro}[htbp]
\caption{Como a ação trata os riscos ambientais e sociais}
\label{qua:socioambiental}
\footnotesize
\renewcommand{\arraystretch}{1.25}
\begin{tabularx}{\textwidth}{|P{3.3cm}|L|P{4.2cm}|}
\hline
\cab{Tema} & \cab{Como a ação atende} & \cab{Evidência} \\ \hline
Avaliação e controle dos impactos &
  Estudos ambientais e sociais; outorga e licenças antes das obras; normas
  técnicas de perfuração & Outorgas e licenças emitidas; relatórios de
  conformidade semestrais \\ \hline
Condições de trabalho &
  Protocolos de saúde e segurança; EPI obrigatório; treinamento das equipes
  das empreiteiras & Relatórios de supervisão e de fiscalização \\ \hline
Prevenção da poluição e uso da água &
  Soluções de esgoto dimensionadas por domicílio; testes de estanqueidade;
  manejo do lodo; gestão de resíduos, ruído e poeira nas obras; hidrômetros &
  Laudos e testes de estanqueidade; relatório técnico aprovado \\ \hline
Saúde e segurança da comunidade &
  Sinalização e proteção de valas; análise da água antes da operação;
  capacitação em uso e manutenção & Análise físico-química; listas de
  presença e certificados \\ \hline
Terra e uso do solo &
  Verificação prévia da titularidade das áreas de poço e reservatório;
  anuências e termos de servidão & Documentação fundiária; anuências
  assinadas \\ \hline
Grupos vulneráveis e comunidades tradicionais &
  Busca ativa em áreas remotas; consultas e métodos participativos; critérios
  de enquadramento sociais & Cadastro das famílias; relatório do trabalho
  técnico-social \\ \hline
Participação e queixas &
  Reuniões comunitárias em linguagem simples; canal de queixas divulgado nas
  comunidades & Registro das queixas e dos tempos de resposta, a cada
  semestre \\ \hline
\end{tabularx}
\fonte{Elaborado pelo IAT (2026) com base em \citeonline[Quadros 25 e 26]{mop2026}.}
\end{quadro}
